\scanpage{69}
\mspage{\handwritten{- 4 -}}

p.37, li.6 : ``composé'' au lieu de ``composée''.

\begin{center}
Commentaires Appendice
\end{center}

p.3 : détailler preuve de 7.4 ? ($Y$ plat sur $S$, ... il en résulte aisément ..., on vérifie aisément ...)

p.1 : dans l'énoncé 7.1, de quel $K^{\bullet}$ s'agit-il ?

p.5, li.5 : réf. pour \handwritten{``}$X_t$ et $X_s$ réguliers $\handwritten{\Rightarrow}$\handwritten{\unsure{``|''}}$X$ régulier\handwritten{''}.

p.6, li.6 : peut-on voir l'argument ?

p.7, li.8 du bas : ``Ce résultat généralise 7.5'' (une fois qu'on a remarqué que, pour $X = S$, les flèches du diagramme de 7.3 sont les flèches identiques de $\underline{Z}$.)

p.7, li.5 du bas : ``récrire'' et non ``réécrire''.

p.8, li.12 du bas : ``horizontales'' et non ``verticales''.

p.8, li.6 du bas : ``du prisme'' et non ``du carré''.

p.9 : nommer les flèches verticales d\deleted{a}e (7.10.2 bis).

p.9, li.6 : ajouter ``(cf. 3.3)'' après \unsure{```|``}implique''.

% Detached marks retain uncertain authorship; small inline boxes preserve their relative placement.
p.9 : références pour $\mathrm{Gr}^{\bullet}_{\mathrm{top}}$ et l'assertion de compatibilité aux filtrations ?\hspace{0.4em}\raisebox{0.35ex}[0pt][0pt]{\parbox[c]{2mm}{\figure{\unsure{\begin{tikzpicture}[x=1mm,y=1mm,baseline=0pt]
\fill (0,0) .. controls (-0.12,0.12) and (-0.09,0.42) .. (0.02,0.52)
 .. controls (0.15,0.63) and (0.30,0.43) .. (0.28,0.24)
 .. controls (0.28,0.06) and (0.14,-0.03) .. (0,0) -- cycle;
\end{tikzpicture}}}}}

Donner aussi le diagramme analogue à (7.10.1 bis) avec \makebox[0pt][l]{\hspace{0.70em}\raisebox{-1.45ex}[0pt][0pt]{\parbox[c]{2mm}{\figure{\unsure{\begin{tikzpicture}[x=1mm,y=1mm,baseline=0pt]
\fill (0,0) .. controls (-0.12,0.10) and (-0.17,0.35) .. (0.01,0.42)
 .. controls (0.16,0.48) and (0.36,0.30) .. (0.25,0.14)
 .. controls (0.17,0.03) and (0.07,-0.05) .. (0,0) -- cycle;
\fill (0.60,1.13) ellipse[x radius=0.34,y radius=0.17,rotate=45];
\end{tikzpicture}}}}}}$\mathrm{Gr}^{\bullet}(\ )_{\underline{Q}}$ pour $g$ projectif, d'intersection complète, et de dimension relative virtuelle donnée.

p.10, li.11 : ``standard'' et non``standart''.

p.10, li.13 : ne faut-il pas ajouter ``connexe'' après ``schéma algébrique'' et ``courbe algébrique'' ? Y a-t-il une référence à l'intention du lecteur qui aimerait voir une démonstration d\deleted{d}e l'assertion en question?

% archive sheet number: 68
\clearpage
